\chapter{Hemoglobin A1c (HbA1c)}

\textbf{Interpretive caution:} HbA1c is useful when the relationship between hemoglobin glycation and average glucose is stable. Results may be misleading when red-cell survival, hemoglobin structure, or the assay method is affected.

\section*{What Is This Test?}
The hemoglobin A1c (HbA1c) test measures the percentage of hemoglobin that has been irreversibly glycated --- the fraction of the A1c subform (hemoglobin A1c) in which glucose is non-enzymatically attached to the N-terminal valine of the beta-globin chain. Because red blood cells circulate for approximately 120 days and glucose attaches slowly and continuously, HbA1c reflects average glycemia over roughly the preceding 8--12 weeks, weighted toward the most recent weeks. HbA1c is used for diagnosis and long-term monitoring of diabetes without requiring fasting. Each 1\% increase in HbA1c corresponds to an estimated average glucose increase of approximately 28--29 mg/dL, although this relationship varies among individuals. The value is associated with diabetes complication risk at the population level but should not be treated as a standalone predictor for an individual patient.

\section*{Alternative Names}
A1c; Glycated Hemoglobin; Glycohemoglobin; Glycosylated Hemoglobin; Hemoglobin A1c; HbA1C; Glycated Hemoglobin (GHb); Total Glycated Hemoglobin (distinct from glycated albumin or fructosamine)

\section*{Principle}
HbA1c forms through a slow, non-enzymatic condensation between glucose and the N-terminal amino group of the beta-globin valine, a reaction that proceeds throughout the red cell's lifespan \cite{tietz2023textbook}. Laboratory methods include cation-exchange HPLC, immunoassay, capillary electrophoresis, and enzymatic methods. Results are standardized to the National Glycohemoglobin Standardization Program (NGSP) against the Diabetes Control and Complications Trial (DCCT) reference method and reported as a percentage; the International Federation of Clinical Chemistry (IFCC) reference method reports mmol/mol. The conversions are: NGSP \% = 0.09148 $\times$ IFCC (mmol/mol) + 2.152, and IFCC (mmol/mol) = 10.929 $\times$ NGSP \% - 23.5 \cite{little2011status}. Red-cell lifespan, hemoglobin variants, recent transfusion, and assay interference can distort the relationship between HbA1c and average glucose.

\section*{History}
The glycated hemoglobin was discovered by Samuel Rahbar in 1968, who noticed an unusual fast-moving hemoglobin fraction on electrophoresis in patients with diabetes \cite{rahbar1968abnormal}. In 1971 Trivelli and colleagues confirmed that this fraction was elevated in diabetes and represented non-enzymatic glycosylation of hemoglobin. The 1976 work of Koenig and Cerami established that the amount of glycosylated hemoglobin reflected integrated blood glucose over the prior weeks, giving rise to the concept of a ``memory'' of glycemia. The landmark Diabetes Control and Complications Trial (DCCT, published 1993) \cite{dcct1993effect} and the United Kingdom Prospective Diabetes Study (UKPDS, published 1998) \cite{ukpds1998intensive} demonstrated that lowering HbA1c reduces microvascular complications (retinopathy, nephropathy, neuropathy) and led to the universal adoption of HbA1c targets and standardized reporting. In 2002 the NGSP harmonized assay standardization worldwide. In 2010 the American Diabetes Association (ADA) formally added HbA1c $\geq$6.5\% to fasting glucose and OGTT as a diagnostic criterion for diabetes \cite{ada2023standards}, a position endorsed by the World Health Organization in 2011 \cite{who2011glycated}.

\section*{Why This Test Is Ordered}
\begin{itemize}
  \item \textbf{Screening for prediabetes and diabetes in asymptomatic adults:} Screening should begin at age 35 years for most adults, and earlier in adults with overweight or obesity plus additional risk factors; repeat testing depends on the initial result and risk profile
  \item \textbf{Diagnosis of diabetes mellitus:} HbA1c $\geq$6.5\% on two separate occasions confirms the diagnosis (a single value $\geq$6.5\% with symptoms of hyperglycemia may suffice)
  \item \textbf{Monitoring glycemic control in established diabetes:} Recommended every 3 months when treatment is being adjusted or glucose is uncontrolled, and every 6 months in stable patients meeting targets
  \item \textbf{Assessment of overall glycemic control in patients on insulin or oral agents} to guide therapy intensification
  \item \textbf{Evaluation during pregnancy in women with pre-existing diabetes} (not for gestational diabetes screening, where it is less sensitive)
  \item \textbf{Prediction of complication risk:} Higher HbA1c correlates with higher risk of retinopathy, nephropathy, and neuropathy
\end{itemize}

\section*{Primary vs Secondary Test}
HbA1c is a primary test for the diagnosis and longitudinal monitoring of diabetes mellitus. When an asymptomatic result is $\geq$6.5\%, it must be repeated on a second sample (with rare exceptions) for confirmation. When used for monitoring, it is the primary outcome measure of glycemic control but is always interpreted together with self-monitored blood glucose values and, in some settings, with glycated albumin (fructosamine).

\section*{How This Test Is Performed}
A venous blood sample is collected in a lavender-top (EDTA) tube, or a small capillary sample may be taken by fingerstick for point-of-care devices (which require NGSP certification and correlation with laboratory methods). Approximately 2--3 mL of whole blood is sufficient. The sample is analyzed on an automated HPLC, immunoassay, capillary electrophoresis, or enzymatic platform. Fasting is not required because HbA1c is unaffected by short-term glucose excursions. Point-of-care A1c testing provides results within minutes during the office visit; laboratory testing typically returns results within 1--2 days. The venipuncture itself takes less than 5 minutes.

\section*{What Preparation Is Needed}
No fasting or special preparation is required. The test can be performed at any time of day. The patient should inform the healthcare provider of all medications (especially erythropoiesis-stimulating agents, iron supplements, and any recent blood transfusion), known hemoglobinopathies, chronic kidney disease, pregnancy, and any history of anemia or recent blood loss, because these conditions alter red blood cell turnover and can make HbA1c unreliable.

\section*{Contraindications and Precautions}
There are no absolute contraindications to the blood draw. HbA1c may be unreliable when red-cell survival is abnormal or a hemoglobin variant interferes with the specific assay. Hemolysis, recent substantial blood loss, recent transfusion, and treatment that rapidly increases young red cells may falsely lower HbA1c. Iron deficiency and vitamin B12 or folate deficiency may falsely raise it in some patients. Hemoglobin variants produce method-dependent effects; the laboratory should select a suitable assay or use glucose-based criteria when the result is discordant with measured glucose. Chronic kidney disease, pregnancy, and sickle-cell trait do not automatically invalidate HbA1c, but require clinical and method-specific interpretation. Fructosamine, glycated albumin, or glucose monitoring may be useful alternatives when HbA1c cannot be interpreted reliably.

\section*{Risks}
Risks are those of routine venipuncture: mild pain, bruising, small hematoma at the collection site, lightheadedness or vasovagal syncope, and extremely rarely, infection or nerve injury. No clinically significant risk is associated with the assay itself.

\section*{What Happens After the Test}
Point-of-care results are discussed at the same visit; laboratory results are usually available within 1--2 days. A diagnostic result should be confirmed when required by the clinical situation, and treatment is individualized according to diabetes type, symptoms, comorbidities, pregnancy status, hypoglycemia risk, and patient preferences. Prediabetes generally prompts lifestyle and risk-factor intervention with follow-up testing at an interval based on risk. For established diabetes, HbA1c is compared with an individualized goal; less than 7\% is appropriate for many nonpregnant adults, but stricter or less stringent goals may be appropriate. HbA1c should be interpreted alongside glucose measurements and, when available, continuous glucose monitoring data.

\section*{Understanding Results}

\subsection*{Below Normal (Less than 5.7\%)}
\begin{itemize}
  \item \textbf{Below the prediabetes threshold:} HbA1c $<$5.7\% is below the diagnostic threshold for prediabetes, but does not by itself establish normal glucose tolerance or exclude dysglycemia when symptoms or risk factors are present
  \item \textbf{Hemolytic anemia or recent blood loss:} Falsely low HbA1c due to shortened red cell survival
  \item \textbf{Recent blood transfusion:} Donor red cells dilute the patient's glycated fraction
  \item \textbf{Pregnancy:} Physiological hemodilution and reduced red cell lifespan lower HbA1c in the second and third trimesters
  \item \textbf{Erythropoiesis-stimulating agents or recent iron therapy:} Young red cells have had less time to accumulate glycation
\end{itemize}

\subsection*{Above Normal (5.7\% and Higher)}
\begin{itemize}
  \item \textbf{Prediabetes (5.7--6.4\%):} Increased risk of progression to diabetes; intensive lifestyle intervention reduces progression by approximately 58\%
  \item \textbf{Diabetes mellitus ($\geq$6.5\%):} Established diabetes; requires confirmation on repeat testing when asymptomatic
  \item \textbf{Poorly controlled diabetes on monitoring:} HbA1c above the individualized target indicates a need to intensify therapy; a change of 1\% corresponds to an average glucose change of approximately 28--29 mg/dL
  \item \textbf{Iron deficiency anemia, B12 or folate deficiency:} May falsely elevate HbA1c in some patients because of altered red-cell turnover; correlate with glucose and correct the underlying disorder
  \item \textbf{Chronic kidney disease with uremia:} Carbamylated hemoglobin may interfere with some methods, falsely elevating the result
\end{itemize}

\section*{Normal Results}
\begin{itemize}
  \item \textbf{Normal (no diabetes):} $<$5.7\% (39 mmol/mol or less in IFCC units) \cite{fischbach2023manual}
  \item \textbf{Prediabetes:} 5.7--6.4\% (39--46 mmol/mol)
  \item \textbf{Diabetes:} $\geq$6.5\% (48 mmol/mol or higher)
  \item \textbf{Treatment target (most nonpregnant adults):} $<$7.0\% (53 mmol/mol)
  \item \textbf{More stringent target (selected younger patients):} $<$6.5\% if achievable without significant hypoglycemia
  \item \textbf{Less stringent target (frail or complex patients):} $<$8.0\% (64 mmol/mol)
\end{itemize}

\section*{Follow-up Tests}
\begin{itemize}
  \item \textbf{Repeat HbA1c:} For confirmation of a diagnostic result, and every 3--6 months for monitoring
  \item \textbf{Fasting plasma glucose:} In conjunction with diagnosis and when discordance with HbA1c requires resolution
  \item \textbf{Oral glucose tolerance test (OGTT):} When HbA1c is borderline or discordant, or for gestational diabetes screening
  \item \textbf{Urine albumin-to-creatinine ratio and serum creatinine:} Screening for diabetic kidney disease
  \item \textbf{Lipid panel:} Cardiovascular risk assessment, performed at diagnosis and periodically
  \item \textbf{Liver function tests:} Baseline before statin therapy and for monitoring of certain glucose-lowering medications
  \item \textbf{Fructosamine or glycated albumin:} When HbA1c is unreliable (hemoglobinopathy, hemolysis, pregnancy, recent transfusion, advanced kidney disease)
  \item \textbf{Continuous glucose monitoring:} To provide the glucose management indicator (GMI) and detailed glucose patterns in insulin-treated patients
\end{itemize}

\section*{References}
\bibliographystyle{unsrtnat}
\bibliography{references}

\clearpage
\section*{Regulatory Status and Authorization Date}
Regulatory status applies to the specific assay, instrument, manufacturer, and intended use rather than to the test name alone. No single approval date can be assigned to this test category. For a commercial assay, verify the current product in the relevant regulator's database: the U.S. Food and Drug Administration (FDA) clearance, approval, or authorization; India's Central Drugs Standard Control Organisation (CDSCO) licence or permission; the European Union In Vitro Diagnostic Regulation (EU IVDR) CE certificate issued through the applicable conformity-assessment route; or the United Kingdom Medicines and Healthcare products Regulatory Agency (MHRA) registration and UKCA/CE status. Laboratory-developed tests may not have an individual FDA, CDSCO, EU IVDR, or MHRA approval date.
\clearpage
